% Ad Familiares 11.1 --- headnote
% ad-familiares-11-01, 16 March 44 BC, Rome

\textit{Decimus Brutus to Marcus Brutus and Cassius, from Rome
on 16 March 44 BC --- Perseus dateline \textit{Scr. Romae xvii
K. Apr. a. 710 (44)}. The letter is written the day after the
Ides, on the morning after Decimus has had Hirtius (the
consul-designate and Caesarian moderate) at his house overnight.
Marcus Brutus and Cassius, the praetors at the head of the
conspiracy, are with their party on the Capitol; Decimus, who
was Caesar's intended governor for Cisalpine Gaul, is sounding
out the ground from his own house in the city.}

\textit{The position laid out here is grim and fast-moving.
Antony, as relayed through Hirtius, will neither hand over the
province nor guarantee the conspirators' safety in Rome, and
the veterans and the urban plebs are stirred up. Decimus
proposes the workaround for which the technical name was
\textit{legatio libera} --- a free, fictitious embassy that
would let the praetors leave Italy without resigning their
offices --- and floats Rhodes as a refuge if that fails. The
calculus in section 4 (Sextus Pompeius in Spain, Caecilius
Bassus in Syria as the only ground to stand on) and the bleak
ladder of section 3 (return to Rome / live in exile / take up
the last resort) show how completely the assassins, within a
day, had lost the initiative.}
